PHYSICS 13 - PROBLEM SOLVING

Collision followed by sliding or a spring: use energy before the collision, momentum across it and energy or work after it. Kinetic energy doesn't survive a sticky collision. Problems 11 , 12 , 13 , 14

"Minimum height," "just makes it" or "leaves the surface": set the normal force to zero at the critical point, then use energy conservation from the start. Problems 8 , 10

"Just starts to move" or "without slipping": static friction is at its maximum, \(\mu N\), and points against the motion that's about to happen. Problems 4 , 5

Find the constraint before writing forces: \(a_1 = 2a_2\) for a movable pulley, \(a = R\alpha\) for a rope on a wheel. Problems 6 , 15

Ropes on a pulley with inertia: Newton's second law for each hanging mass, \(\sum\tau = I\alpha\) for the pulley and \(a = R\alpha\) to tie them together. Keep one sign convention throughout. Problems 15 , 17 , 18

Is it SHM? Show that the net force is \(-(\text{constant})\times x\); then \(\omega = \sqrt{\text{constant}/m_{\text{eff}}}\), where a rolling object's \(m_{\text{eff}}\) includes its spin. Problems 21 , 22

\(pV\)-diagram sign tables: get \(W_{\text{by}}\) from whether \(V\) grows, \(\Delta U\) from whether \(pV\) (and so \(T\)) grows, then \(Q = \Delta U + W_{\text{by}}\). Check whether the question asks for work done on the gas or by it. Problems 26 , 27 , 28

Check the limits: free fall makes \(g_{\text{eff}} = 0\); a horizontal string would need infinite \(\omega\); on a composite pulley the heavier mass wins only if its torque \(MR\) is bigger. Problems 15 , 16 , 20
